\documentclass[11pt]{article}
\usepackage[margin=1in]{geometry}
\usepackage{amsmath,amssymb,amsthm,booktabs,tabularx,longtable}
\usepackage[T1]{fontenc}
\usepackage{lmodern}
\usepackage[hidelinks]{hyperref}
\newtheorem{theorem}{Theorem}
\newtheorem{lemma}[theorem]{Lemma}
\newtheorem{corollary}[theorem]{Corollary}
\newcommand{\coeff}[1]{[x^{#1}]}
\title{Prime-index divisibility in Hanna's OEIS A381355\\
\large Revision 13, COMMENTS, and shifted-power siblings}
\author{Haobo Ma\thanks{ChronoAI Pte Ltd, Singapore.}\and
Wenlin Zhang\thanks{National University of Singapore; ChronoAI Pte Ltd, Singapore.}}
\date{November 6, 2026}
\begin{document}
\maketitle
\begin{abstract}
We prove Paul D. Hanna's conjecture in the COMMENTS of OEIS A381355,
revision 13 (March 12, 2025): its $n$th term is divisible by the $n$th
prime for every $n>1$. The source series of A381353 exists uniquely over
$\mathbb Q$ and has integer coefficients. More generally, if positive
integers $m_n$ satisfy $\gcd(m_n,n!)=1$, the equations
$\coeff n F^{m_n}=0$ for $n\ge2$, with $F(0)=1$ and prescribed integer
$\coeff1F$, define an integral series whose logarithmic coefficient
$\coeff{n-1}F'/F$ is divisible by $m_n$.
The same argument proves companion integrality and the prime divisibility
conjecture in A393867. A parity argument proves its oddness conjecture and
the even-run conjecture in A393868. We also correct the coefficient index
in A393866, FORMULA (2), and give a general nonzero-diagonal congruence.
\end{abstract}

\section{The conjecture and conventions}
Let $p_n$ be the $n$th prime, with $p_1=2$. Hanna's A381353
\cite{source} defines a generating function by
\begin{equation}\label{eq:source}
 F(x)=1+x+\sum_{n\ge2}f_nx^n,\qquad \coeff n F(x)^{p_n}=0\quad(n\ge2).
\end{equation}
The initial coefficients $f_0=f_1=1$ are fixed by its PROGRAM section.
The sequence A381355 \cite{target} is defined, with offset one, by
\begin{equation}\label{eq:target}
 \frac{xF'(x)}{F(x)}=\sum_{n\ge1}a_nx^n,
 \qquad a_n=\coeff{n-1}\frac{F'(x)}{F(x)}.
\end{equation}
Its COMMENTS section states: ``Conjecture: a(n) is divisible by prime(n)
for n > 1.'' We establish it and the existence and integrality implicit
in these definitions. All series below are formal; no analytic convergence
is required.

The argument is a coefficient construction followed by a formal derivative
identity. Primality can be replaced by coprimality with a factorial.
We give that general statement so that the arithmetic hypothesis and the
companion integrality property follow from the same calculation.
The equivalent A381353, revision 19, COMMENTS, statement for
$\coeff{n-1}F'/F$ says ``is divisible by prime(n) for n > 1.''
The observation first appears in
\href{https://oeis.org/history/view?seq=A381353&v=8}{A381353, revision 8, COMMENTS}
(March 11, 2025, 21:29:48 EDT), before the companion entry A381355.
The two entries cite each other and give numerical tables rather than proofs.
Section~\ref{sec:shift} treats the later family A393866--A393868, posed
March 29, 2026, with its two logarithmic conjectures and a parity-run remark.

\section{A coefficient congruence}
\begin{lemma}\label{lem:low}
Let $N\ge1$ and $m\ge1$ be integers with $\gcd(m,N!)=1$.
For every $B\in1+x\mathbb Z[[x]]$,
\[
 B(x)^m\equiv1\pmod{(m,x^{N+1})}.
\]
Equivalently, $m$ divides $\coeff j B^m$ for every $1\le j\le N$.
\end{lemma}
\begin{proof}
For $1\le k\le N$, we have $\gcd(k,m)=1$ and
$k\binom mk=m\binom{m-1}{k-1}$. Thus $m\mid\binom mk$,
where $\binom mk=0$ if $k>m$.
Write $B=1+U$ with $U\in x\mathbb Z[[x]]$ and expand
$(1+U)^m=\sum_{k=0}^m\binom mkU^k$.
Terms with $k>N$ have degree greater than $N$; every remaining nonconstant
term has coefficients divisible by $m$.
\end{proof}

\section{Construction and divisibility}
\begin{theorem}\label{thm:general}
Fix $c\in\mathbb Z$ and a sequence of positive integers $(m_n)_{n\ge2}$
such that $\gcd(m_n,n!)=1$ for every $n\ge2$.
There is exactly one $F\in\mathbb Q[[x]]$ such that
\[
 \coeff0F=1,\qquad\coeff1F=c,\qquad
 \coeff nF^{m_n}=0\quad(n\ge2).
\]
It belongs to $\mathbb Z[[x]]$. If $a_n=\coeff{n-1}F'/F$, then
$a_n\in\mathbb Z$ and $m_n\mid a_n$ for every $n\ge2$.
\end{theorem}
\begin{proof}
Suppose $f_0=1,f_1=c,f_2,\ldots,f_{n-1}$ have been chosen and put
$B_n=\sum_{j=0}^{n-1}f_jx^j$. For any $t\in\mathbb Q$, the binomial
expansion gives, since $2n>n$,
\begin{equation}\label{eq:triangular}
 \coeff n(B_n+tx^n)^{m_n}=\coeff nB_n^{m_n}+m_nt.
\end{equation}
Consequently the only possible choice is
\begin{equation}\label{eq:recurrence}
 f_n=-\frac{\coeff nB_n^{m_n}}{m_n}.
\end{equation}
By induction $B_n$ has integer coefficients; Lemma~\ref{lem:low} with
$N=n$ shows that the quotient in \eqref{eq:recurrence} is an integer.
Choosing these coefficients for every $n$ defines $F$. Later coefficients
cannot change an earlier equation, because its coefficient of degree $n$
depends only on $f_0,\ldots,f_n$. This proves existence over $\mathbb Z$.
Equation~\eqref{eq:triangular}, inductively over $\mathbb Q$, proves
uniqueness there as well.

Since $F(0)=1$, its inverse belongs to $\mathbb Z[[x]]$: if
$F^{-1}=\sum g_jx^j$, then $g_0=1$ and
$g_n=-\sum_{j=1}^nf_jg_{n-j}$. Hence $L=F'/F$ has integer coefficients.
Fix $n\ge2$ and set $m=m_n$. Formal differentiation gives
\begin{equation}\label{eq:derivative}
 (F^m)'=mLF^m.
\end{equation}
The coefficient of degree $n-1$ on the left is $n\coeff nF^m=0$.
As $m$ is a nonzero integer, \eqref{eq:derivative} implies
\[
 \coeff{n-1}(LF^m)=0.
\]
Lemma~\ref{lem:low} shows $F^m\equiv1\pmod{(m,x^{n+1})}$.
Multiplication by the integral series $L$ therefore yields
$\coeff{n-1}(LF^m)\equiv\coeff{n-1}L=a_n\pmod m$.
It follows that $m\mid a_n$.
\end{proof}

\begin{corollary}[Hanna's conjecture]\label{cor:hanna}
The series in \eqref{eq:source} exists uniquely over $\mathbb Q$, has
integer coefficients, and its sequence \eqref{eq:target} satisfies
$p_n\mid a_n$ for every $n>1$.
\end{corollary}
\begin{proof}
The $n$th prime satisfies $p_n>n$: among the integers from $2$ through
$n$ there are at most $n-1$ primes. Thus $p_n$ is coprime to $n!$.
Apply Theorem~\ref{thm:general} with $c=1$ and $m_n=p_n$.
\end{proof}

\begin{corollary}[The companion coefficients]\label{cor:sibling}
For the same $F$, every $\coeff jF^{p_n}/p_n$ is an integer when
$n\ge2$ and $1\le j<p_n$. In particular,
\[
 b_n=\frac{\coeff{n-1}F^{p_n}}{p_n}\in\mathbb Z\quad(n\ge2).
\]
The logarithmic reconstruction is
\[
 F(x)=\exp\left(\sum_{n\ge1}\frac{a_n}{n}x^n\right)
 \quad\text{in }\mathbb Q[[x]].
\]
\end{corollary}
\begin{proof}
Apply Lemma~\ref{lem:low} with $m=p_n$ and $N=p_n-1$; integrality of
$F$ was proved above. For the last identity, let the series in parentheses
be $S(x)$. Equation~\eqref{eq:target} gives $S'=F'/F$.
Thus $(F\exp(-S))'=0$. Over $\mathbb Q$, a formal series with derivative
zero is constant; its constant term here is one.
\end{proof}

The restriction $n>1$ is necessary: $a_1=1$ is not divisible by $p_1=2$.
The general theorem also applies to any prime schedule $q_n>n$, to any
integer $c$, and to the composite schedule $m_n=q_n^2$ with $q_n>n$ prime.
The coprimality hypothesis cannot simply be dropped: $c=1,m_2=2$
forces $f_2=-1/2$.
These last two unrestricted readings are not asserted in the source:
they fail at $n=1$ and at $(c,m_2)=(1,2)$, respectively.

\begin{corollary}[Prescribed diagonals]\label{cor:diagonal}
Under the exponent hypotheses of Theorem~\ref{thm:general}, let
$c,d_2,d_3,\ldots$ be any integers. The equations
$F(0)=1$, $\coeff1F=c$, and
$\coeff nF^{m_n}=m_nd_n$ for $n\ge2$ determine a unique rational
series. It is integral, and
\[
 \coeff{n-1}F'/F\equiv nd_n\pmod {m_n}\quad(n\ge2).
\]
\end{corollary}
\begin{proof}
At step $n$, replace \eqref{eq:recurrence} by
$f_n=d_n-\coeff n B_n^{m_n}/m_n$. The same lemma proves integrality,
and \eqref{eq:triangular} proves uniqueness. Differentiation now gives
$\coeff{n-1}(LF^{m_n})=nd_n$. Since $F^{m_n}\equiv1$ through
degree $n-1$ modulo $m_n$, this reduces to the claimed congruence.
\end{proof}

\section{The shifted-power family}\label{sec:shift}
Write $U=\sum_{j\ge0}u_jx^j$ for the series in A393866, revision 17,
NAME and PROGRAM \cite{shift}, defined by
\begin{equation}\label{eq:shift}
 u_0=1,\qquad
 \coeff nU^{p_n}=p_n\coeff{n-1}U^{p_n}\quad(n\ge1).
\end{equation}
Its $n=1$ equation gives $u_1=1$. Put
$\alpha_n=\coeff{n-1}U'/U$; these are the offset-one terms of
A393867 \cite{shiftlog}. Its COMMENTS contains two conjectures:
``a(n) is divisible by prime(n) for n > 1'' and ``all terms are odd.''

\begin{corollary}\label{cor:shift}
Equation~\eqref{eq:shift} has a unique rational solution, which is integral,
and $p_n\mid\alpha_n$ for every $n>1$.
More generally, for the exponents of Theorem~\ref{thm:general}, any
integers $c,t_2,t_3,\ldots$ define a unique integral series $V$ by
\[
 V(0)=1,\quad\coeff1V=c,\quad
 \coeff nV^{m_n}=m_nt_n\coeff{n-1}V^{m_n}\quad(n\ge2),
\]
and $m_n\mid\coeff{n-1}V'/V$ for every $n\ge2$.
\end{corollary}
\begin{proof}
The coefficient of degree $n-1$ on the right depends only on the
already determined coefficients. Thus the triangular construction applies
with $d_n=t_n\coeff{n-1}B_n^{m_n}$. It is integral and divisible by
$m_n$ by Lemma~\ref{lem:low}. The proof of
Corollary~\ref{cor:diagonal}, degree by degree, gives the conclusion.
Take $c=t_n=1$ and $m_n=p_n$ for \eqref{eq:shift}.
\end{proof}

\begin{corollary}\label{cor:odd}
Every $\alpha_n$ is odd for $n\ge1$. Moreover,
$u_{2k}\equiv u_{2k+1}\pmod2$ for every $k\ge0$.
Consequently every finite maximal parity run in A393866 has even length.
This is A393868, revision 9, COMMENTS: ``It appears that all terms
are even'' \cite{runs}.
\end{corollary}
\begin{proof}
Work in $\mathbb F_2[[x]]$, retaining the notation $U,u_j$ for their
reductions. The first pair is $u_0=u_1=1$. Suppose the pairs with
indices below $2k$ agree, where $k\ge1$. Set $n=2k+1$ and
$H=U^{(p_n-1)/2}=\sum_{r\ge0}h_rx^r$. Here $p_n$ is odd,
$H(0)=1$, and $U^{p_n}=UH^2$. Reducing \eqref{eq:shift} gives
\[
 0=\coeff{2k+1}(1+x)UH^2
   =\sum_{j=0}^k(u_{2j+1}+u_{2j})h_{k-j}^{\,2}
   =u_{2k+1}+u_{2k}.
\]
Induction proves the pairing. Hence $U=(1+x)G(x^2)$ for
$G=\sum_{k\ge0}u_{2k}x^k$ with $G(0)=1$. Formal differentiation
in characteristic two yields $U'=G(x^2)$ and
$U'/U=(1+x)^{-1}=\sum_{j\ge0}x^j$. This proves oddness of
every $\alpha_n$. A parity change can occur only between two successive
pairs, so every completed run has even length.
\end{proof}

\paragraph{A printed index and an unresolved limit.}
A393866, revision 17, FORMULA (2), instead uses
$\coeff n U'/U\equiv0\pmod{p_n}$ for $n>1$.
This printed conjecture is false at $n=2$: from
$U=1+x+2x^2+10x^3+\cdots$ one obtains
$U'/U=1+3x+25x^2+\cdots$, and $3\nmid25$.
The corrected exponent is $n-1$, exactly as in
Corollary~\ref{cor:shift} and the indexed PARI calculation in A393866.
The separate A393866 COMMENTS observation of a limiting ratio to the
primorial near $0.078$ is not settled by these arguments.

\section{Coverage and verification}
{
\small
\begin{longtable}{@{}>{\raggedright\arraybackslash}p{0.40\linewidth}>{\raggedright\arraybackslash}p{0.55\linewidth}@{}}
\toprule
Source item & Mathematical coverage \\
\midrule
\endfirsthead
\toprule
Source item & Mathematical coverage (continued) \\
\midrule
\endhead
A381355, rev.\ 13, COMMENTS & Conjecture proved in Corollary~\ref{cor:hanna}. \\
A381353, rev.\ 19, NAME and PROGRAM & Not posed; integral existence and rational uniqueness established in Theorem~\ref{thm:general}. \\
A381353, COMMENTS & Same divisibility statement, proved in Corollary~\ref{cor:hanna}. \\
A381353, rev.\ 19, EXAMPLE, RELATED SEQUENCES & Not posed; all-index integrality of $b_n$ established in Corollary~\ref{cor:sibling}. \\
A381355, rev.\ 13, EXAMPLE, logarithmic series & Proved identity, Corollary~\ref{cor:sibling}. \\
A393866, rev.\ 17, NAME and PROGRAM & Not posed; integral existence and rational uniqueness established in Corollary~\ref{cor:shift}. \\
A393867, rev.\ 16, COMMENTS, first conjecture & Proved: prime divisibility, Corollary~\ref{cor:shift}. \\
A393867, rev.\ 16, COMMENTS, second conjecture & Proved: all terms odd, Corollary~\ref{cor:odd}. \\
A393868, rev.\ 9, COMMENTS & Proved: all finite parity-run lengths even, Corollary~\ref{cor:odd}. \\
A393866, rev.\ 17, FORMULA (2) & Refuted-and-corrected: $n=2$ gives $25\not\equiv0\pmod3$; correct degree $n-1$. \\
A393866, rev.\ 17, COMMENTS, primorial limit & Not covered: existence and value of the suggested limit. \\
Prescribed diagonals and composite exponents & Not posed; proved in Corollaries~\ref{cor:diagonal} and~\ref{cor:shift}. \\
\bottomrule
\end{longtable}
}
The displayed power tables illustrate the defining equations.

All three formal sources below are fixed at revision\\
\texttt{de81f09d1b5868b1ef5f62c9c8cab94962481c39}.\\
The pinned \href{https://github.com/the-omega-institute/trureturing/blob/de81f09d1b5868b1ef5f62c9c8cab94962481c39/D5/S1/Recurrence/Invariants/HannaPrimeIndexLogDivisibility.lean}{\texttt{HannaPrimeIndexLogDivisibility.lean}}
contains the declaration
\begin{quote}\small
\path{D5.S1.Recurrence.Invariants.HannaPrimeIndexLogDivisibility.result}.
\end{quote}
The declaration \texttt{result} proves \texttt{claim}, which asserts
existence of a rational source
series and, for every such series and every $n>1$, an integer $z$ with
$\coeff{n-1}(F'/F)=p_nz$. This is the existence and prime-divisibility
part of Corollary~\ref{cor:hanna}. Theorem~\ref{thm:general}, its
coprime-composite extension and Corollary~\ref{cor:sibling} have the
written proofs above.

For the shifted series, the pinned
\href{https://github.com/the-omega-institute/trureturing/blob/de81f09d1b5868b1ef5f62c9c8cab94962481c39/D5/S1/Recurrence/Parity/PrimePowerShiftLogDerivative.lean}{\texttt{PrimePowerShiftLogDerivative.lean}}
proves the defining equations and uniqueness among integral
series in \texttt{generating\_equation} and \texttt{generating\_unique};
\texttt{hanna\_conjecture} proves $p_n\mid\alpha_n$ for $n>1$, and
\texttt{printed\_formula\_false} proves failure of the printed degree-$n$
formula. The pinned
\href{https://github.com/the-omega-institute/trureturing/blob/de81f09d1b5868b1ef5f62c9c8cab94962481c39/D5/S1/Recurrence/Parity/PrimePowerShiftLogDerivativeOdd.lean}{\texttt{PrimePowerShiftLogDerivativeOdd.lean}}
proves coefficient pairing in \texttt{generating\_coeff\_pair}
and oddness of $\alpha_n$ for $n\ge1$ in \texttt{a393867\_odd}.
The finite-run consequence and the general prescribed-diagonal and
variable-shift results have the written proofs above.

Exact finite checks in \texttt{verify.py} cover the original sequence
through degree $64$ and $24$ combinations of exponent schedule and integer
initial slope through degree $24$. The all-index assertions follow from
the written proofs above. The companion file \texttt{verify\_siblings.py}
checks the shifted series and logarithmic coefficients through degree $64$,
all $32$ coefficient pairs available through degree $63$, completed parity
runs, the printed-formula example, and $24$ prescribed-diagonal families
through degree $24$. These are finite checks rather than proofs of an
infinite statement.

\small
\begin{thebibliography}{9}
\raggedright
\bibitem{target}
Paul D. Hanna, \emph{A381355}, The On-Line Encyclopedia of Integer
Sequences, posed March 11, 2025, revision 13, March 12, 2025,
NAME, COMMENTS and EXAMPLE.
\url{https://oeis.org/A381355};
\href{https://oeis.org/history/view?seq=A381355&v=13}{revision 13}.

\bibitem{source}
Paul D. Hanna, \emph{A381353}, The On-Line Encyclopedia of Integer
Sequences, posed March 11, 2025, revision 19, January 1, 2026,
NAME, COMMENTS, EXAMPLE and PROGRAM.
\url{https://oeis.org/A381353};
\href{https://oeis.org/history/view?seq=A381353&v=19}{revision 19}.

\bibitem{shift}
Paul D. Hanna, \emph{A393866}, The On-Line Encyclopedia of Integer
Sequences, March 29, 2026, revision 17, NAME, COMMENTS, FORMULA,
EXAMPLE and PROGRAM. \url{https://oeis.org/A393866};
\href{https://oeis.org/history/view?seq=A393866&v=17}{revision 17}.

\bibitem{shiftlog}
Paul D. Hanna, \emph{A393867}, The On-Line Encyclopedia of Integer
Sequences, posed March 29, 2026, revision 16, August 20, 2026,
COMMENTS and EXAMPLE. \url{https://oeis.org/A393867};
\href{https://oeis.org/history/view?seq=A393867&v=16}{revision 16}.

\bibitem{runs}
Paul D. Hanna, \emph{A393868}, The On-Line Encyclopedia of Integer
Sequences, March 29, 2026, revision 9, COMMENTS and EXAMPLE.
\url{https://oeis.org/A393868};
\href{https://oeis.org/history/view?seq=A393868&v=9}{revision 9}.

\end{thebibliography}

\noindent\small Copyright 2026 Haobo Ma and Wenlin Zhang.
This note is licensed under Creative Commons Attribution 4.0 International
(CC BY 4.0): \url{https://creativecommons.org/licenses/by/4.0/}.
\end{document}
